\documentclass[11pt,a4paper]{article}

% ---------------------------------------------------------
% PACKAGES
% ---------------------------------------------------------
\usepackage[utf8]{inputenc}
\usepackage[T1]{fontenc}
\usepackage{lmodern}
\usepackage{amsmath, amssymb, amsfonts}
\usepackage{bm}
\usepackage{graphicx}
\usepackage{hyperref}
\usepackage{authblk}
\usepackage{geometry}
\usepackage{cite}
\usepackage{physics}
\usepackage{tensor}
\usepackage{enumitem}

\geometry{margin=1in}

% ---------------------------------------------------------
% TITLE
% ---------------------------------------------------------
\title{\textbf{Companion XIII: The Neutrino and Dark Matter Sectors in the Relaxation-Driven Cyclic Cosmology (RDCC)}}

\author[1]{Michael Lehmann \\ \texttt{mi.lehmann@gmx.de}}
\affil[1]{\small Independent Researcher, Relaxation-Driven Cyclic Cosmology (RDCC) Framework}

\date{\small \today}

% ---------------------------------------------------------
% DOCUMENT START
% ---------------------------------------------------------
\begin{document}

\maketitle

\begin{abstract}
The Relaxation-Driven Cyclic Cosmology (RDCC) provides a unified framework in which the large-scale evolution of the Universe, the emergence of cosmic asymmetries, and the sectoral structure of matter arise from a single dynamical principle: relaxation across CPT-conjugate sectors. In this Companion, we develop a combined treatment of the neutrino and dark matter sectors, which together constitute the non-baryonic matter content of RDCC. We show that both components inherit their properties from the same sectoral thermodynamics, the same relaxation-induced mass suppression mechanism, and the same inter-sector gravitational coupling. Neutrinos and dark matter jointly shape the growth of structure, the evolution of perturbations, and the observational signatures in the CMB, lensing, and large-scale structure. This Companion establishes the unified sectoral matter picture of RDCC and provides the foundation for future inference and forecasting analyses.
\end{abstract}

% ---------------------------------------------------------
% SECTION 1 — INTRODUCTION
% ---------------------------------------------------------
\section{Introduction}

The Relaxation-Driven Cyclic Cosmology (RDCC) provides a unified dynamical framework in which the large-scale evolution of the Universe, the emergence of cosmic asymmetries, and the sectoral structure of matter arise from a single principle: relaxation across CPT-conjugate sectors. In this picture, the matter content of the Universe is naturally divided into two sectors related by a global CPT transformation, each evolving with opposite thermodynamic arrows and exchanging information only through gravity and the relaxation dynamics.

The neutrino and dark matter sectors form the non-baryonic matter content of RDCC. Both components inherit their effective properties from the same sectoral thermodynamics, the same relaxation-induced mass suppression mechanism, and the same inter-sector gravitational coupling. This Companion develops a unified treatment of these two sectors, showing that they form a dynamically coupled subsystem whose properties encode direct information about the relaxation parameter $\alpha$, the inter-sector entropy flow, and the global cyclic boundary conditions.

The structure of this Companion is as follows. Section 2 develops the neutrino sector, including the mass hierarchy, Dirac/Majorana structure, and relaxation-induced suppression. Section 3 introduces the dark matter sector, which in RDCC arises as baryonic matter in the CPT-conjugate sector. Section 4 develops the joint sectoral thermodynamics, while Section 5 derives the relaxation-induced mass suppression mechanism. Section 6 presents the cosmological signatures, and Section 7 provides forecasts for Euclid and future surveys.

% ---------------------------------------------------------
% SECTION 2 — THE NEUTRINO SECTOR IN RDCC
% ---------------------------------------------------------
\section{The Neutrino Sector in RDCC}

Neutrinos play a central role in the RDCC framework. Their small masses, long free-streaming lengths, and weak interactions make them sensitive probes of the relaxation dynamics and the inter-sector thermodynamic structure. In RDCC, neutrinos inherit their properties from three ingredients:

\begin{enumerate}[label=(\alph*)]
    \item the global CPT structure of the two sectors,
    \item the relaxation-induced suppression of effective masses,
    \item the sectoral entropy flow across the cyclic boundary.
\end{enumerate}

\subsection{Mass hierarchy and sectoral symmetry}

In RDCC, the neutrino mass hierarchy is not an arbitrary parameter but emerges from the interplay between relaxation and sectoral thermodynamics. The CPT-conjugate sector contains a mirror neutrino population with opposite thermodynamic arrow, and the two populations are related by a sectoral symmetry that constrains the allowed mass patterns. The effective neutrino masses in our sector take the form
\begin{equation}
m_{\nu,i}^{\text{eff}} = m_{\nu,i}^{(0)} \, e^{-\alpha \, S_{\text{rel}}},
\end{equation}
where $m_{\nu,i}^{(0)}$ are the bare masses, $\alpha$ is the relaxation parameter, and $S_{\text{rel}}$ is the sectoral entropy difference.

\subsection{Dirac/Majorana structure}

The RDCC framework allows both Dirac and Majorana neutrinos, but the relaxation dynamics naturally suppresses Majorana mass terms relative to Dirac terms. This leads to a hybrid structure in which the effective neutrino masses are predominantly Dirac-like, with small Majorana components induced by inter-sector relaxation. The resulting phenomenology differs from both standard Dirac and standard Majorana scenarios.

\subsection{Relaxation-induced mass suppression}

The relaxation dynamics suppresses neutrino masses through the exponential factor $e^{-\alpha S_{\text{rel}}}$. This mechanism provides a natural explanation for the smallness of neutrino masses without requiring fine-tuning or new high-energy scales. The suppression is tied to the inter-sector entropy flow and therefore encodes information about the global cyclic boundary conditions.

\subsection{Impact on structure formation}

Neutrinos affect structure formation through their free-streaming, which suppresses the growth of perturbations on small scales. In RDCC, this suppression is modified by the relaxation dynamics, leading to a scale-dependent signature that differs from the standard $\Lambda$CDM prediction. The combined neutrino–dark-matter sector produces a characteristic imprint on the matter power spectrum, lensing potential, and CMB anisotropies.
% ---------------------------------------------------------
% SECTION 3 — THE DARK MATTER SECTOR IN RDCC
% ---------------------------------------------------------
\section{The Dark Matter Sector in RDCC}

In the RDCC framework, dark matter (DM) arises not as a new particle species but as baryonic matter in the CPT-conjugate sector. This identification follows directly from the global CPT structure of the two-sector Universe and the relaxation dynamics that govern their evolution. The dark matter sector is therefore not exotic: it is composed of ordinary baryons, photons, and neutrinos in the opposite thermodynamic arrow, whose gravitational influence is felt in our sector.

\subsection{Sectoral identification of dark matter}

The CPT-conjugate sector contains a baryonic plasma whose thermodynamic arrow runs opposite to ours. Because the two sectors interact only gravitationally (and through relaxation), the baryons in the conjugate sector behave as cold, collisionless matter from the perspective of our sector. This provides a natural explanation for the cold and pressureless nature of dark matter without invoking new particle species.

\subsection{Effective equation of state}

The effective equation of state of dark matter in RDCC is determined by the thermodynamic history of the conjugate sector. As the conjugate sector cools along its opposite arrow, its baryonic matter becomes effectively cold and pressureless when viewed from our sector. The resulting equation of state is
\begin{equation}
w_{\text{DM}}^{\text{eff}} \approx 0,
\end{equation}
consistent with the observational requirements for cold dark matter.

\subsection{Free-streaming and clustering}

Unlike standard CDM, the dark matter in RDCC inherits a small but non-zero free-streaming scale from the thermodynamic history of the conjugate sector. This leads to a mild suppression of small-scale structure, which is further modified by the relaxation dynamics. The combined effect produces a characteristic scale-dependent signature in the matter power spectrum.

\subsection{Gravitational coupling between sectors}

The two sectors interact only through gravity and the relaxation dynamics. The gravitational potential felt in our sector receives contributions from both sectors:
\begin{equation}
\Phi = \Phi_{\text{ours}} + \Phi_{\text{CPT}},
\end{equation}
where $\Phi_{\text{CPT}}$ is sourced by the baryonic matter in the conjugate sector. This coupling leads to distinctive signatures in lensing, structure formation, and the CMB.

% ---------------------------------------------------------
% SECTION 4 — JOINT SECTORAL THERMODYNAMICS
% ---------------------------------------------------------
\section{Joint Sectoral Thermodynamics}

The neutrino and dark matter sectors share a common thermodynamic origin in the RDCC framework. Both components arise from the interplay between the global CPT structure, the relaxation dynamics, and the inter-sector entropy flow. This section develops the joint thermodynamic picture that governs the evolution of the two sectors.

\subsection{Sectoral entropy flow}

The two sectors exchange entropy across the cyclic boundary, leading to a sectoral entropy difference $S_{\text{rel}}$ that drives the relaxation dynamics. The entropy flow determines the effective masses of neutrinos and dark matter, as well as their thermodynamic histories.

\subsection{Temperature evolution}

The temperature evolution of the two sectors is governed by their opposite thermodynamic arrows. In our sector, the temperature decreases with time, while in the conjugate sector it increases. This leads to a temperature asymmetry that affects the effective properties of neutrinos and dark matter.

\subsection{Impact on cosmological observables}

The joint thermodynamics of the two sectors affects a wide range of cosmological observables, including the matter power spectrum, the lensing potential, and the CMB anisotropies. The combined neutrino–dark-matter sector produces a characteristic imprint that differs from the standard $\Lambda$CDM prediction.


\section{Thermal Mirror Sound Speed Before Mirror Recombination}
\label{sec:mirror-sound-speed}

In the RDCC framework, the dark matter component is identified with the
baryonic matter of the CPT-conjugate mirror sector. As shown in earlier
sections of this Companion, the mirror sector is thermodynamically colder,
$T_{\mathrm{mirror}} < T_{\mathrm{visible}}$, and undergoes its own recombination
epoch at an earlier scale factor. Before this mirror recombination, the
mirror baryons form a tightly coupled photon--baryon fluid with a finite,
temperature-dependent sound speed. Consequently, the dark matter component
cannot be treated as a perfectly cold, pressureless fluid during this epoch.

\subsection{Mirror Photon--Baryon Fluid}

Prior to mirror recombination, the mirror baryons and mirror photons are
tightly coupled through Thomson scattering. The sound speed of this fluid is
given by the standard expression
\begin{equation}
c_{s,\mathrm{mirror}}^{2}(a)
=
\frac{1}{3}
\left[
1
+
\frac{3\rho_{b,\mathrm{mirror}}(a)}{4\rho_{\gamma,\mathrm{mirror}}(a)}
\right]^{-1},
\label{eq:mirror-sound-speed}
\end{equation}
where $\rho_{b,\mathrm{mirror}}$ and $\rho_{\gamma,\mathrm{mirror}}$ denote the mirror
baryon and mirror photon energy densities, respectively. Because the mirror
sector is colder, the ratio $\rho_{b,\mathrm{mirror}}/\rho_{\gamma,\mathrm{mirror}}$
grows earlier than in the visible sector, and the sound speed decreases more
rapidly with scale factor.

\subsection{Transition at Mirror Recombination}

Mirror recombination occurs when the mirror temperature drops below the
hydrogen binding threshold,
\begin{equation}
T_{\mathrm{mirror}}(a_{\mathrm{rec,mirror}})
\simeq 0.3\,\mathrm{eV},
\end{equation}
which corresponds to a scale factor
\begin{equation}
a_{\mathrm{rec,mirror}}
=
\frac{T_{\mathrm{visible}}}{T_{\mathrm{mirror}}}
a_{\mathrm{rec,visible}}
\ll
a_{\mathrm{rec,visible}}.
\end{equation}
After this epoch, the mirror photons decouple, the pressure support collapses,
and the mirror baryons behave as a pressureless fluid. The sound speed
therefore transitions sharply from Eq.~\eqref{eq:mirror-sound-speed} to
\begin{equation}
c_{s,\mathrm{mirror}}^{2}(a>a_{\mathrm{rec,mirror}}) \approx 0.
\end{equation}

\subsection{Impact on the Matter Power Spectrum}

The finite sound speed before mirror recombination suppresses the growth of
perturbations on scales below the mirror-sector sound horizon. The comoving
sound horizon is
\begin{equation}
r_{s,\mathrm{mirror}}(a)
=
\int_{0}^{a}
\frac{c_{s,\mathrm{mirror}}(a')}{a'^{2}H(a')}
\,\mathrm{d}a'.
\end{equation}
Modes with $k r_{s,\mathrm{mirror}} \gtrsim 1$ experience acoustic oscillations
and partial suppression, analogous to the baryon acoustic oscillations in the
visible sector but shifted to earlier times and smaller scales.

This mechanism produces a smooth, continuous suppression of the matter power
spectrum at large wavenumbers,
\begin{equation}
P(k)
\;\longrightarrow\;
P(k)\,
\mathcal{S}_{\mathrm{mirror}}(k),
\end{equation}
where the suppression factor is approximately
\begin{equation}
\mathcal{S}_{\mathrm{mirror}}(k)
\simeq
\left[
1
+
\left(\frac{k}{k_{\mathrm{mirror}}}\right)^{2}
\right]^{-\alpha},
\qquad
k_{\mathrm{mirror}}
\sim
r_{s,\mathrm{mirror}}^{-1}.
\end{equation}
The exponent $\alpha$ is of order unity and depends on the detailed thermal
history of the mirror sector.

\subsection{Implications for RDCC and CLASS Implementation}

This finite sound speed is a structural prediction of RDCC and provides a
clean physical origin for the smooth high-$k$ suppression of the matter power
spectrum. Unlike warm dark matter models, where the suppression arises from
free-streaming, the RDCC suppression is acoustic in nature and tied directly
to the mirror-sector thermodynamics.

The corresponding modification of the perturbation equations is implemented
in Companion~XVI, where the mirror dark matter is treated as a fluid with
sound speed $c_{s,\mathrm{mirror}}^{2}(a)$ until mirror recombination. This
ensures that the CLASS evolution accurately reflects the thermal history of
the mirror sector and preserves the predictive structure of RDCC.


% ---------------------------------------------------------
% SECTION 5 — RELAXATION-INDUCED MASS SUPPRESSION
% ---------------------------------------------------------
\section{Relaxation-Induced Mass Suppression}

A key feature of the RDCC framework is the relaxation-induced suppression of effective masses. This mechanism provides a natural explanation for the smallness of neutrino masses and the cold nature of dark matter.

\subsection{Exponential suppression mechanism}

The effective masses of neutrinos and dark matter are suppressed by an exponential factor of the form
\begin{equation}
m_{\text{eff}} = m_{0} \, e^{-\alpha \, S_{\text{rel}}},
\end{equation}
where $m_{0}$ is the bare mass, $\alpha$ is the relaxation parameter, and $S_{\text{rel}}$ is the sectoral entropy difference.

\subsection{Implications for neutrinos}

The exponential suppression naturally explains the smallness of neutrino masses without requiring fine-tuning or new high-energy scales. The suppression is tied to the inter-sector entropy flow and therefore encodes information about the global cyclic boundary conditions.

\subsection{Implications for dark matter}

The dark matter in RDCC inherits its effective mass from the baryonic matter in the conjugate sector. The relaxation-induced suppression ensures that the dark matter behaves as cold, collisionless matter, consistent with observational requirements.

\subsection{Combined signatures}

The combined neutrino–dark-matter sector produces a characteristic scale-dependent signature in the matter power spectrum, lensing potential, and CMB anisotropies. These signatures provide a powerful probe of the relaxation dynamics and the sectoral thermodynamics.
% ---------------------------------------------------------
% SECTION 6 — COSMOLOGICAL SIGNATURES
% ---------------------------------------------------------
\section{Cosmological Signatures}

The combined neutrino–dark-matter sector of RDCC produces a rich set of cosmological signatures that differ from the predictions of the standard $\Lambda$CDM model. These signatures arise from the relaxation-induced mass suppression, the sectoral thermodynamics, and the gravitational coupling between the two sectors.

\subsection{Matter power spectrum}

The matter power spectrum receives contributions from both sectors. Neutrino free-streaming suppresses power on small scales, while the dark matter inherits a mild free-streaming scale from the thermodynamic history of the conjugate sector. The combined effect produces a characteristic scale-dependent signature that can be used to constrain the relaxation parameter $\alpha$.

\subsection{Weak lensing}

Weak lensing is particularly sensitive to the combined neutrino–dark-matter sector. The lensing potential receives contributions from both sectors, and the relaxation dynamics modifies the growth of structure in a way that produces distinctive signatures in the lensing power spectrum, bispectrum, and peak statistics.

\subsection{CMB anisotropies}

The CMB anisotropies are affected by the combined neutrino–dark-matter sector through the integrated Sachs–Wolfe effect, the lensing potential, and the early-time suppression of perturbations. The relaxation dynamics introduces a scale-dependent modification that differs from the standard $\Lambda$CDM prediction.

\subsection{Large-scale structure}

The large-scale structure of the Universe provides a powerful probe of the combined neutrino–dark-matter sector. The relaxation dynamics modifies the growth rate of perturbations, leading to a scale-dependent signature in the redshift-space distortions and the velocity divergence field.

% ---------------------------------------------------------
% SECTION 7 — FORECASTS AND ALPHA-SENSITIVITY
% ---------------------------------------------------------
\section{Forecasts and \texorpdfstring{$\alpha$}{alpha}-Sensitivity}

The combined neutrino–dark-matter sector provides one of the most sensitive probes of the relaxation parameter $\alpha$. This section presents forecasts for Euclid and future surveys, showing how the joint sectoral matter content can be used to constrain the relaxation dynamics.

\subsection{Fisher forecasts}

Fisher forecasts show that the combined neutrino–dark-matter sector provides strong constraints on $\alpha$, with sensitivities that exceed those of standard cosmological probes. The scale-dependent signatures in the matter power spectrum, lensing potential, and CMB anisotropies provide complementary information that can be used to break degeneracies.

\subsection{Multi-probe constraints}

Combining multiple probes, including weak lensing, galaxy clustering, and CMB lensing, provides even stronger constraints on $\alpha$. The joint analysis reveals a characteristic pattern of scale-dependent signatures that is unique to the RDCC framework.

\subsection{Future surveys}

Future surveys, including Euclid, LSST, and CMB-S4, will provide high-precision measurements of the matter power spectrum, lensing potential, and CMB anisotropies. These measurements will allow for precise constraints on the relaxation parameter $\alpha$ and the sectoral thermodynamics.

% ---------------------------------------------------------
% CONCLUSION
% ---------------------------------------------------------
\section{Conclusion}

This Companion has developed a unified treatment of the neutrino and dark matter sectors in the Relaxation-Driven Cyclic Cosmology (RDCC). Both components arise from the interplay between the global CPT structure, the relaxation dynamics, and the inter-sector thermodynamics. The combined neutrino–dark-matter sector produces a rich set of cosmological signatures that differ from the predictions of the standard $\Lambda$CDM model.

The relaxation-induced mass suppression provides a natural explanation for the smallness of neutrino masses and the cold nature of dark matter. The joint sectoral thermodynamics and the gravitational coupling between the two sectors produce distinctive signatures in the matter power spectrum, lensing potential, and CMB anisotropies.

The combined neutrino–dark-matter sector provides one of the most sensitive probes of the relaxation parameter $\alpha$, and future surveys will allow for precise constraints on the relaxation dynamics and the sectoral thermodynamics. This Companion establishes the unified sectoral matter picture of RDCC and provides the foundation for future inference and forecasting analyses.

% ---------------------------------------------------------
% REFERENCES
% ---------------------------------------------------------
\begin{thebibliography}{99}

\bibitem{RDCC-Flagship}
Michael Lehmann, ``Relaxation-Driven Cyclic Cosmology,'' RDCC Flagship Paper (2026).

\bibitem{RDCC-Companion-IV}
Michael Lehmann, ``Companion IV: The Global CPT State and the Mirror Sector in RDCC,'' (2026).

\bibitem{RDCC-Companion-XII}
Michael Lehmann, ``Companion XII: Sectoral Baryon Asymmetry from Opposite Thermodynamic Arrows in RDCC,'' (2026).

\bibitem{RDCC-Companion-XVI}
Michael Lehmann, ``Companion XVI: The RDCC Boltzmann Code,'' (2026).

\bibitem{RDCC-Companion-XXVI}
Michael Lehmann, ``Companion XXVI: Gravitational Lensing in RDCC,'' (2026).

\bibitem{RDCC-Companion-LXXXIII}
Michael Lehmann, ``Companion LXXXIII: Fisher Forecasts for $\alpha$ Using Topological Transport,'' (2026).

\end{thebibliography}

\end{document}
